% Figure: the function zoo plotted over a shared horizontal range
% so the reader can compare scaling behaviours on the same axes. A
% linear function, a quadratic, an exponential, a logarithm, and a
% sine all share x in [0, 6].
\begin{figure}[htbp]
\centering
\begin{tikzpicture}[scale=1.05,
  ax/.style={->, draw=engblack, thick},
  tick/.style={draw=engblack!60, thin},
  lin/.style={draw=engblack, thick},
  quad/.style={draw=engblue, thick},
  expo/.style={draw=engred, thick},
  loga/.style={draw=engamber, thick},
  trig/.style={draw=enggray, thick, dashed},
  ann/.style={font=\footnotesize, align=left},
]
% --- axes ---
\draw[ax] (0, 0) -- (6.5, 0) node[below right, ann] {$x$};
\draw[ax] (0, 0) -- (0, 5.0) node[left, ann] {$y$};
\foreach \x in {1, 2, 3, 4, 5, 6} {
  \draw[tick] (\x, -0.05) -- (\x, 0.05);
  \node[ann, below, yshift=-2pt] at (\x, 0) {\x};
}
\foreach \y in {1, 2, 3, 4} {
  \draw[tick] (-0.05, \y) -- (0.05, \y);
  \node[ann, left, xshift=-2pt] at (0, \y) {\y};
}

% --- linear: y = 0.5 x ---
\draw[lin, domain=0:6, samples=2, smooth] plot (\x, {0.5*\x});
\node[ann, anchor=south west, draw=engblack, fill=white, inner sep=2pt]
  at (5.2, 2.8) {$y = x/2$};

% --- quadratic: y = 0.12 x^2 ---
\draw[quad, domain=0:6, samples=40, smooth] plot (\x, {0.12*\x*\x});
\node[ann, anchor=south, draw=engblue, fill=white, inner sep=2pt]
  at (5.5, 4.4) {$y = x^{2}/8$};

% --- exponential: y = 0.2 e^{0.6 x} ---
\draw[expo, domain=0:5.5, samples=40, smooth] plot (\x, {0.2*exp(0.6*\x)});
\node[ann, anchor=west, draw=engred, fill=white, inner sep=2pt]
  at (5.5, 4.85) {$y = 0.2 e^{0.6 x}$};

% --- logarithm: y = ln(x + 1) ---
\draw[loga, domain=0:6, samples=40, smooth]
  plot (\x, {ln(\x + 1)});
\node[ann, anchor=south, draw=engamber, fill=white, inner sep=2pt]
  at (5.4, 2.1) {$y = \ln(x + 1)$};

% --- sinusoid: y = 1.5 + sin(x*180/pi) ---
\draw[trig, domain=0:6, samples=60, smooth]
  plot (\x, {1.5 + sin(\x * 57.2958)});
\node[ann, anchor=west, draw=enggray, fill=white, inner sep=2pt]
  at (0.3, 0.5) {$y = 1.5 + \sin x$};

\end{tikzpicture}
\caption[The elementary-function zoo on a shared range]{Five
elementary functions plotted on the same horizontal range $x \in [0,
6]$ so the reader can compare scaling behaviours on the same eyes.
The linear function $y = x/2$ rises by a constant slope. The
quadratic $y = x^{2}/8$ stays below the linear until $x = 4$, then
overtakes it. The exponential $0.2 e^{0.6 x}$ begins below both,
crosses the linear near $x = 3$, and outruns the quadratic by $x =
5$. The logarithm $\ln(x + 1)$ rises rapidly near the origin and
flattens. The sinusoid oscillates between $0.5$ and $2.5$ around
its mean $1.5$. The diagnostic discipline: knowing which of the five
shapes a dataset most resembles is the first move before fitting any
model.}
\label{fig:vol02:ch02:function-zoo}
\end{figure}
